\scanpage{52}
% manuscript line 1
\deleted{Montrons que \illegible{1} $a\in E$} Forme qui peut s'écrire, en \deleted{\unsure{considérant}} \deleted{\unsure{le vecteur}} \inserted{posant} $u=\bar a\otimes b$

% manuscript line 2
\deleted{\unsure{Maintenant on a de plus}}

% manuscript line 3
\[
\deleted{(15)\quad\text{\unsure{On a de plus}}}\qquad (u.x,y)=(u,\bar x\otimes y) \tag*{(15)}
\]

% manuscript line 4
\deleted{\unsure{On a|Comme}\ \illegible{1}} \deleted{\unsure{où dans} le second}
(dans le second mb, on prend \deleted{le} produit scalaire
dans l'espace préhilbertien $\bar E\otimes F$), formule qui
restera valable \inserted{par linéarité} pour \deleted{toute} $u\in\bar E\otimes F$. On en
conclut \deleted{que} $|(ux,y)|\leq\|u\|_2\,\|\bar x\otimes y\|_2=\|u\|_2\|\bar x\|\,\|y\|
=\|u\|_2\,\|x\|\,\|y\|$, d'où résulte que l'opérateur
défini par $u$ a une norme $\leq\|u\|_2$. Par
suite, l'application $u\to\tilde u$ de $\bar E\otimes F$ dans $L(E,F)$
se prolonge par continuité en une application linéaire
continue de \deleted{$\bar E\otimes$} norme $\leq 1$ de $\bar E\overset{(2)}{\otimes}F$. Cette
application est d'ailleurs biunivoque, car si $(e_i)_{i\in I}$ est
une base orthonormale dans $E$ \inserted{(donc aussi dans $\bar E$)}, $(f_j)_{j\in J}$ une base
orthonormale dans $F$, alors $\bar E\overset{(2)}{\otimes}F$ s'identifie
à l'espace \deleted{d\unsure{es} $l^2$} des matrices infinies
$u=(u_{ij})_{ij\in I\times J}$ qui sont de carré sommable, \inserted{(prop. 1)} et
il est \deleted{\unsure{évident}} \inserted{trivial} que l'opérateur \deleted{$E$} défini par $u$
a \deleted{la} \inserted{pour} matrice $(u_{ij})$, donc $u$ est nul si cet
opérateur est nul. \deleted{Nous ne distinguerons pas par la
suite entre les $u\in\bar E\overset{(2)}{\otimes}F$ et les op. qu'ils définissent}

% manuscript line 23
\uline{Définition 2} Soient $E$ et $F$ deux espaces de Hilbert.
On appelle op. de H.-S. de $E$ dans $F$, \deleted{\unsure{toutes}} \inserted{les} applications
linéaires de $E$ dans $F$ définies par des éléments
de $\bar E\overset{(2)}{\otimes} F$. L'esp. de ces op., \deleted{est} muni de la structure

% manuscript line 27
% archive sheet number: 53
